\subsection{Project Overview}\label{project-overview}

In this assignment, our goal was to develop a Monte Carlo simulation
program to explore the conformational landscape of a 500-Carbon linear
polymer chain (polyethylene) with varying torsional angles. The
simulation generates the initial conformation (with \& without
Hydrogens), measures torsional energy rotations, end-to-end distance,
and intramolecular Lennard-Jones interactions, and finally analyzes \&
visualizes the results.

After creating a serial processing version of the program,
parallelization techniques involving both MPI and OpenMP were
incorporated. These were compared to the serial version to demonstrate
both their individual and combined High Performance Computing
capabilities \& benefits.

Many techniques taught in class were crucial in the development of this
program, including: - \textbf{Git/Github} as a code respository \&
teamwork coordination platform - \textbf{Makefile} to automate the
compilation and execution of the code - \textbf{Shell scripting} for
facilitating server deployment - \textbf{MPI \& openMP} to parallelize
the code - \textbf{Python} for post-processing and plotting